\section{Discussion and Outlook}
\label{sec:discussion_outlook}

\paragraph{What this work contributes.}
The central contribution is a physically interpretable duality geometry for mean-force equilibrium:
one reflection exchanges weak and strong influence representations, and a second reflection handles
orientation labeling.
Together they convert difficult asymptotic sectors into tractable partner calculations while keeping
the same analytic family.

\paragraph{Why this is useful beyond one model.}
The framework only needs two structural ingredients:
\begin{enumerate}
  \item a positive crossover coordinate with distinct weak/strong asymptotics, and
  \item a closed representative map on control labels.
\end{enumerate}
Whenever these exist, one can separate physical content (state family) from computational strategy
(representative choice).
This is a practical advantage for low-temperature and strong-coupling studies where direct expansions
are fragile.

\paragraph{Scope and limits.}
The dual maps are statements about representative control, not about parameter-level equality of
physical states.
The spin-boson analysis provides an explicit existence example with closure-level exact formulas,
including spectral-class critical couplings and entropy branch transfer.
These results should be interpreted at the level of the closed mean-force representation used here;
broader universality claims require model-by-model analysis.

\paragraph{Relation to established literature.}
For spin-boson criticality, the benchmark remains nonperturbative universality analyses:
Ohmic KT behavior and super-Ohmic delocalization from the classic dissipative two-state framework,
and sub-Ohmic criticality from NRG/QMC and related approaches
\cite{leggettDynamicsDissipativeTwostate1987,weissQuantumDissipativeSystems2012,bullaQuantumPhaseTransitionSubohmic2003,vojtaQuantumPhaseTransitionsSubohmic2005,winterQuantumPhaseTransitionSubohmic2009,bullaNumericalRenormalizationGroup2008}.
For mean-force thermodynamics, the present duality-geometry perspective complements rather than
replaces existing structural and operational formulations
\cite{gelinThermodynamicsSubensembleCanonical2009a,hiltHamiltonianMeanForce2011,talknerColloquiumStatisticalMechanics2020,rivasStrongCouplingThermodynamics2020,burkeStructureHamiltonianMean2024}.

\paragraph{Experimental and numerical relevance.}
The geometric framing suggests straightforward diagnostics:
map data into \(y=\log\chi\), identify self-dual contour \(\chi=1\), and test partner reconstructions
across branches.
The orientation-gauge affine reflection gives an additional filter for distinguishing radial
invariants from direction-sensitive observables in fitted effective descriptions.

\paragraph{Natural extensions.}
Three directions are immediate:
\begin{enumerate}
  \item multi-channel couplings with several \(\chi_a\), where \((\mathbb{Z}_2)^k\) branch reflections
  may emerge,
  \item finite-\(N\) systems with richer eigenlabel symmetries than qubit \(\mathbb{Z}_2\), and
  \item thermodynamic-cycle settings where representative transfer can reduce control-cost modeling in
  strong-coupling operation windows.
\end{enumerate}

\paragraph{Bottom line.}
Mean-force duality is most useful when treated as a control language: not a hidden symmetry of
parameters, but a map between computationally complementary views of the same dressed equilibrium
state.
